\frfilename{mt23.tex}
\versiondate{17.11.04}
\copyrightdate{1995}


\def\chaptername{Teorema Radon-Nikod\'ym}
\newchapter{23}
\def\chaptername{Teorema Radon-Nikod\'ym}

Dalam Bab 22, saya membahas integral tak tentu dari fungsi-fungsi
terintegralkan pada $\Bbb R$, dan memberikan apa yang saya harap Anda anggap
sebagai uraian yang memuaskan, baik mengenai
fungsi-fungsi yang merupakan integral tak tentu (fungsi-fungsi kontinu
mutlak) maupun mengenai cara menemukan fungsi yang diintegralkan untuk
menghasilkan integral tak tentu tersebut (dengan mendiferensialkannya).
Untuk ruang ukur umum, tidak tersedia struktur yang dapat menghasilkan
perumusan sesederhana itu; meskipun demikian, pertanyaan-pertanyaan yang sama
tetap dapat diajukan dan, sampai taraf tertentu, dijawab.

Bagian pertama bab ini memperkenalkan perangkat dasar yang diperlukan,
yaitu konsep fungsional `aditif terhitung' dan penguraiannya menjadi bagian
positif dan negatif. Teorema utama dibahas dalam bagian kedua: integral tak
tentu adalah fungsional aditif yang `kontinu sejati'; pada ruang
$\sigma$-hingga, fungsional tersebut adalah fungsional aditif terhitung yang
`kontinu mutlak'. Dalam \S233 saya membahas penerapan tunggal terpenting
teorema tersebut, yakni penggunaannya untuk menyediakan konsep `ekspektasi
bersyarat'. Ini merupakan salah satu konsep utama teori probabilitas --
sebagaimana sangat mungkin telah Anda ketahui; tetapi bentuknya di sini
merupakan perumuman yang jauh menjangkau dari konsep elementer probabilitas
bersyarat suatu kejadian jika kejadian lain diketahui, dan memerlukan seluruh
kekuatan teori umum ukuran dan integrasi sebagaimana dikembangkan dalam Jilid
1 dan bab ini. Saya menyertakan beberapa catatan tentang fungsi cembung,
hingga dan termasuk versi-versi ketaksamaan Jensen (233I-233J).

Selagi kita berada dalam ranah teori ukuran `murni', saya menggunakan
kesempatan ini untuk membahas beberapa topik tambahan.
Saya memulai dengan sejumlah konstruksi yang pada dasarnya elementer,
yaitu ukuran citra, jumlah ukuran, dan ukuran integral tak tentu;
menurut saya perinciannya memerlukan sedikit perhatian,
dan saya menguraikannya dalam \S234.
Gagasan yang agak lebih mendalam diperlukan untuk menangani `transformasi
terukur'. Dalam \S235 saya memaparkan teknik yang diperlukan untuk menyediakan
landasan abstrak bagi metode umum integrasi dengan substitusi, disertai
uraian terperinci mengenai syarat-syarat cukup agar rumus berbentuk

\Centerline{$\int g(y)dy=\int g(\phi(x))J(x)dx$}

\noindent berlaku.

\discrpage
